\begin{abstract}
The prevalence of large foundation models has surged, yet the financial and computational costs associated with training such models from the ground up remain immense. As a result, the focus has shifted to adapting pre-trained models for specific downstream applications with minimal resource expenditure. This work explores Orthogonal Finetuning (OFT), a theoretically motivated strategy for finetuning foundation models. Despite its capacity for robust generalization, OFT’s parameter footprint is substantial, a consequence of the dimensionality of orthogonal matrices involved. To enhance parameter efficiency, we analyze OFT from the viewpoint of information transmission and delineate several core requirements for improved efficiency. Drawing inspiration from the Cooley-Tukey fast Fourier transform, which achieves efficient information routing, we introduce a compact orthogonal parameterization via butterfly structures. Integrating this parameterization into OFT leads to a new finetuning technique, termed Orthogonal Butterfly~(BOFT). BOFT generalizes OFT, treating it as a special instance within a broader framework for orthogonal finetuning. To validate our method, we present extensive experimental results, showing effective adaptation of large vision transformers, language models, and text-to-image diffusion models across a spectrum of downstream vision and language tasks.
\end{abstract}

\section{Introduction}

Notable foundation models such as ChatGPT~\cite{brown2020language,chen2021evaluating} and Stable Diffusion~\cite{rombach2022high} have highlighted the extraordinary generalization capabilities of large-scale neural architectures. The continuous improvement in their performance has been accompanied by exponential growth in parameter counts

(\emph{e.g.}, the 175B parameters in GPT-3). This escalation in model size renders end-to-end training increasingly unattainable for most practitioners. Consequently, significant advancements hinge on methods that allow practitioners to repurpose existing foundation models for new tasks efficiently. Specifically, it is crucial to develop strategies for rapid adaptation to downstream tasks without the cost of retraining entire models. Efficient adaptation is typically achieved via three major approaches: \emph{model finetuning}~\cite{hulora2022,zaken2022bitfit,guo2020parameter,raffel2020exploring,qiu2023controlling,zhang2022adaptive,chavan2023one}, which adjusts only a select subset of the parameters within the original model architecture; \emph{adapter tuning}~\cite{rebuffi2017learning,houlsby2019parameter,pfeiffer2020adapterfusion,lin2020exploring,he2021towards}, which augments the architecture with additional trainable modules; and \emph{prompt tuning}~\cite{li2021prefix,lester2021power}, which appends trainable prefix tokens to model inputs. Among these, model finetuning is notable for its straightforwardness, its strong performance, and the important property that it does not incur additional latency at inference time.

Model finetuning fundamentally aims to maintain the similarity between the pretrained and finetuned models according to specific metrics, thereby safeguarding the knowledge acquired during pretraining. As an example, prevailing finetuning strategies often utilize a modest learning rate. This pragmatic choice limits divergence between the initial and adapted models. Recognizing the inherent structure of weight matrices, more systematic methods focus on conserving their relational properties, particularly the pairwise angular relationships between neurons. This perspective motivates the development of the Orthogonal Finetuning~(OFT) framework~\cite{qiu2023controlling}, in which all neurons within a single layer are collectively transformed by an identical orthogonal matrix. This design ensures the pairwise angles between neurons are rigorously maintained during finetuning. While OFT achieves strong generalization and efficient convergence in settings such as text-to-image diffusion models~\cite{qiu2023controlling}, it faces practical limitations because orthogonal matrices in high-dimensional spaces require managing a large number of parameters. To enhance parameter efficiency, OFT incorporates a block-diagonal design, substantially reducing the parameter count. Nonetheless, this efficiency is accompanied by trade-offs—namely, the imposed block-diagonal sparsity enforces a rigid structure, with independent orthogonal transformations within each block. Such a sparsity constraint, while advantageous in terms of parameter reduction, potentially brings unintended inductive biases (\emph{e.g.}, diminishing representational capacity and limiting the ability to emulate conventional linear transformations). Furthermore, optimal determination of the sparsity pattern for these blocks remains an unresolved question.

A central challenge in addressing this issue lies in constructing a dense orthogonal matrix while maintaining a parameter-efficient design. Although it might initially appear impractical—since a dense orthogonal matrix of size $d$ inherently requires $\mathcal{O}(d^2)$ parameters—we introduce a novel method to synthesize such matrices via products of several sparse matrices. Our method leverages the observation that computational time can be increased to save on storage, as the orthogonal matrix is encoded through successive multiplication of sparse factors, requiring repeated computations during training. To contextualize the matrix decomposition, we cast the formation of a dense orthogonal matrix as a task of information transmission. Within this viewpoint, constructing a dense orthogonal matrix by composing sparse matrices becomes analogous to passing information through a grid-structured network. Adopting this transmission-based interpretation allows for a rich variety of sparse factorization strategies that effectively constrain the trainable parameter count while preserving the ability to approximate dense transformations.

For parameter-efficient realization in this framework, we turn to the butterfly architecture, inspired by the Cooley-Tukey fast Fourier transform algorithm~\cite{cooley1965algorithm}, which facilitates highly efficient exchanges of information across nodes~\cite{klasing1994broadcasting}. The butterfly pattern in the Cooley-Tukey algorithm naturally gives rise to a corresponding sparse matrix product decomposition, referred to as butterfly factorization. Using this factorization, we can build a $d$-dimensional dense matrix as a product of $\mathcal{O}(\log d)$ sparse factors, where each sparse matrix has only $\mathcal{O}(d)$ non-zero elements. Enforcing orthogonality at each factor results in an efficient orthogonal parameterization requiring merely $\mathcal{O}(d\log d)$ parameters—a substantial reduction compared to the full parameter set. Building upon this compact orthogonal representation, we introduce Orthogonal Butterfly~(BOFT), a novel parameter-efficient finetuning technique. BOFT generalizes and extends OFT, serving as a flexible framework for orthogonal finetuning. Both the block-diagonal and butterfly decompositions rely on inducing sparsity within orthogonal matrices, which serves to decrease the total number of trainable parameters. It is instructive to compare this approach to the low-rank parametrization employed by LoRA~\cite{hulora2022}; both low-rank and sparse matrices are prominent categories within structured matrices~\cite{candes2011robust}, central to achieving parameter efficiency.

In contrast to OFT’s use of a block-diagonal design to balance expressivity with regularization, BOFT leverages the butterfly pattern to achieve a more gradual interpolation between matrices drawn from the full orthogonal group (for expressiveness) and the identity matrix (for regularity). This approach enables the identification of a more compact hypothesis space within the orthogonal group for applications in downstream tasks. Notably, as the butterfly structure underpins several efficient linear operations—including the discrete Fourier and discrete cosine transforms—we posit that our structurally constrained strategy promotes parameter efficiency by introducing an implicit inductive bias, thereby enhancing generalizability and mitigating overfitting. This insight is rooted in the observation that the butterfly structure readily reproduces numerous well-known linear transforms, a capability absent from the block-diagonal framework employed in OFT. The principal contributions of our work are as follows:

\begin{itemize}[leftmargin=*,nosep]
    \item We approach the parameter efficiency challenge in orthogonal finetuning by introducing a new information transmission paradigm, which reformulates the creation of a parameter-efficient dense orthogonal matrix as a problem of information propagation over a graph with a grid-like structure.
    \item Building on the butterfly framework exemplified by the Cooley-Tukey algorithm, we present Orthogonal Butterfly—a new, parameter-efficient orthogonal finetuning method—within this information transmission setting.
    \item We offer theoretical motivations that elucidate how BOFT can drastically minimize the quantity of trainable parameters, while still maintaining ample expressivity and robust training dynamics. The factorization structure inherent to BOFT further allows for compelling weight interpolation capabilities (see Figure~\ref{fig:boft_interpolation}).
    \item For the first time, we extend orthogonal finetuning to a spectrum of tasks beyond controllable text-to-image generation~\cite{qiu2023controlling}, thereby highlighting its broad potential as a general-purpose model adaptation technique. We demonstrate BOFT's effectiveness on multiple adaptation benchmarks, spanning domains from computer vision to natural language processing. Our method achieves notable improvements over current state-of-the-art baselines, underscoring its exceptional parameter efficiency and capacity for generalization.
\end{itemize}

\section{Related Work}

\textbf{Parameter-efficient finetuning (PEFT).} As the size and capabilities of foundation models (\emph{e.g.}, \cite{brown2020language,radford2021learning,rombach2022high,kirillov2023segment}) continue to increase, there is heightened interest in adapting these models for specific downstream tasks while minimizing the number of trainable parameters~\cite{treviso2023efficient,ding2023parameter,lialin2023scaling}. A wide range of PEFT approaches have been explored~\cite{houlsby2019parameter,aghajanyan2020intrinsic,hulora2022,edalati2022krona,wang2022adamix,gheini2021cross,zaken2022bitfit,guo2020parameter,sung2021training,ansell2022composable,lester2021power,li2021prefix,vu2022spot,he2021towards,mao2021unipelt,karimi2021compacter,liu2022few,sung2022lst,chen2022parameter,jia2022visual,chen2022adaptformer,zhang2022neural,jie2023fact,lian2022scaling,luo2023towards}; among these, reparameterization-based strategies~\cite{aghajanyan2020intrinsic,hulora2022,edalati2022krona} are most pertinent to our study. LoRA~\cite{hulora2022}, for instance, finetunes pretrained models by augmenting weight matrices with the product of two learnable low-rank matrices, demonstrating notable results on language processing tasks. While LoRA uses a fixed rank for the added matrices, more recent works~\cite{zhang2022adaptive,valipour2022dylora,zhang2023increlora} introduce techniques for adaptive rank assignment across different network layers, thereby optimizing the distribution of the parameter budget. Owing to its straightforward implementation, low-rank reparameterization techniques have become widely adopted~\cite{dettmers2023qlora,zi2023delta,chavan2023one}. Building on the role of hyperspherical energy in shaping generalization~\cite{liu2018learning,liu2021orthogonal}, \cite{qiu2023controlling} present orthogonal finetuning (OFT), a distinct strategy for adjusting text-to-image diffusion models. Here, OFT involves learning an orthogonal transformation applied within the same layer, leading to enhanced generalization properties and more stable training behavior when compared to LoRA. However, OFT typically entails a larger number of learnable parameters than LoRA, presenting an opportunity to enhance its efficiency. Furthermore, the utility of OFT outside the scope of text-to-image diffusion model adaptation~\cite{qiu2023controlling} remains an open question. BOFT addresses the parameter efficiency challenge of OFT by introducing butterfly factorization. As a result, the efficacy of orthogonal finetuning can now be systematically extended to general model adaptation scenarios.

\textbf{Butterfly structures.} The radix-2 Cooley-Tukey algorithm recursively decomposes the $N$-point discrete Fourier transform into two subproblems of size $\frac{N}{2}$, naturally leading to a butterfly structure that can be expressed as the product of several sparse matrices, collectively known as a butterfly matrix. These matrices have been leveraged in various applications, including orthogonal matrix parameterization to bypass pivoting during Gaussian elimination and enhance computational efficiency~\cite{parker1995random}, promoting stable recurrent neural network training~\cite{jing2017tunable}, and enabling efficient kernel approximation~\cite{munkhoeva2018quadrature}. Recent research~\cite{dao2019learning,dao2019kaleidoscope} has explored learning rapid linear mappings via butterfly parametrizations, while others~\cite{chen2022pixelated,dao2022monarch} have applied these matrices to accelerate neural network training. Beyond these, butterfly structures are instrumental in fast matrix-vector product computation~\cite{michielssen1996multilevel,o2010algorithm}, compressing large matrices~\cite{li2015butterfly}, and efficient network communication~\cite{klasing1994broadcasting,li2003linear,rajasingh2016transmission}. Distinct from previous applications, our work is centered on employing butterfly structures to boost the parameter efficiency of OFT.

\section{An Information Transmission View on Orthogonal Finetuning}

We begin by outlining the foundational concepts of OFT. In order to adapt a pretrained weight matrix, OFT reformulates the updated weight matrix as the product of a trainable orthogonal matrix and the fixed pretrained weight matrix. Unlike LoRA, which modifies the weight matrix by adding a low-rank trainable matrix, OFT incorporates a multiplicative orthogonal matrix to update the weights. While LoRA achieves parameter-efficiency through the use of a low-rank parameterization, classical OFT attains efficiency by employing a block-diagonal orthogonal matrix structure~\cite{qiu2023controlling}. The efficiency of OFT improves as the size of the diagonal blocks decreases. Figure~\ref{fig:comp_lora} illustrates a visual comparison between the approaches. The underlying rationale for using orthogonal transformations when fine-tuning weight matrices is to maintain the pairwise angular relationships among neurons~\cite{liu2018learning,liu2021orthogonal,qiu2023controlling}, thus retaining much of the semantic information acquired during pretraining. Specifically, OFT involves learning an orthogonal matrix $\bm{R}\in\mathbb{R}^{d\times d}$ associated with a pretrained linear layer $\bm{W}^0\in\mathbb{R}^{d\times n}$, altering the forward computation from $\bm{z}=(\bm{W}^0)^\top\bm{x}$ to $\bm{z}=(\bm{R}\bm{W}^0)^\top\bm{x}$, where $\bm{x}\in\mathbb{R}^{d}$ is the input vector and $\bm{z}\in\mathbb{R}^{n}$ is the output. For initialization that maintains the initial network function, $\bm{R}$ is set to the identity matrix at the start. To preserve the orthogonality constraint on $\bm{R}$ during fine-tuning, we adopt Cayley parameterization, following~\cite{qiu2023controlling,liu2021orthogonal}: $\bm{R}=(\bm{I}+\bm{Q})(\bm{I}-\bm{Q})^{-1}$, where $\bm{Q}$ is a skew-symmetric matrix such that $\bm{Q}=-\bm{Q}^\top$. To further reduce the number of trainable parameters, $\bm{R}$ is constructed as block-diagonal, i.e., $\text{diag}(\bm{R}_1,\bm{R}_2,\cdots,\bm{R}_r)$, where each block $\bm{R}_i\in\mathcal{R}^{b\times b}$ is an orthogonal matrix and $br=d$. However, this parameter-efficient construction comes with a limitation: it presumes that the dimensions of each neuron (i.e., each column vector of $\bm{W}^0$) can be partitioned into $r$ groups, with each group transformed independently by different orthogonal blocks. Although the block-diagonal scheme shows empirical success, there is no principled justification for grouping a neuron's dimensions based solely on their indices, which motivates the exploration of dense orthogonal matrices. This leads to a pertinent question: \emph{Is it possible to realize a dense orthogonal matrix that retains parameter-efficiency?}

To tackle this problem, we suggest constructing a dense orthogonal matrix by multiplying several sparse orthogonal matrices together. To create a consistent and intuitive way to analyze the sparsity structure involved in orthogonal matrix decomposition, we conceptualize the formation of a dense orthogonal matrix within OFT as an instance of information transmission. Concretely, expressing a dense matrix $\bm{R}\in\mathbb{R}^{d\times d}$ as a product of $m$ square matrices, $\thickmuskip=2mu \medmuskip=2mu \bm{R} = \bm{B}_m\bm{B}_{m-1}\cdots\bm{B}_1$, can be interpreted as routing information across a grid composed of $d\times (m+1)$ nodes, as depicted in Figure~\ref{fig:info_trans}. This perspective draws inspiration from the realization that a dense $d \times d$ matrix can be seen as a fully connected mapping from one set of $d$ nodes to another. In the context of $\bm{R}$, an entry $R_{ij}=0$ signifies that information cannot be transferred from node $j$ to node $i$, while a nonzero value signifies possible transmission. As a result, the decomposition of $\bm{R}$ into the product $\bm{B}_m\bm{B}_{m-1}\cdots\bm{B}_1$ can equivalently be seen as a sequence of information propagation steps determined by the graphs defined by each $\bm{B}_i$. The passage of information starts according to the pattern in $\bm{B}_1$ and continues sequentially, culminating with $\bm{B}_n$. For example, as shown in Figure~\ref{fig:info_trans}, the factorization $\bm{R} = \bm{B}_5\bm{B}_4\bm{B}_3\bm{B}_2\bm{B}_1$ is illustrated alongside its respective sparsity configurations and the resulting induced graph. The diagram in Figure~\ref{fig:info_trans} unfolds the process of multiplying these matrices. In the corresponding graph, the matrix $\bm{B}_i$ prescribes the connectivity from nodes at layer $i$ to those at layer $i+1$. More precisely, an element $(j_1,j_2)$ in $\bm{B}_i$ encodes the presence (or absence, if zero) of a directed connection from the $j_2$-th node at the $i$-th layer to the $j_1$-th node at the subsequent $(i+1)$-th layer. For the product $\bm{B}_5\bm{B}_4\bm{B}_3\bm{B}_2\bm{B}_1$ to produce a dense matrix, each initial node at the first layer must be able to reach every terminal node at the sixth layer through some information pathway. If we instead examine $\bm{R}=\bm{B}_4\bm{B}_3\bm{B}_2\bm{B}_1$, connecting first-level (source) nodes to fifth-level (sink) nodes, we observe that information starting at node $1$ is unable to reach node $3$, which reveals that the resulting product has a zero at the $(3,1)$ position in its sparsity pattern. Similar relationships between the product’s induced graph and its sparsity pattern appear for shorter products, such as $\bm{R}=\bm{B}_3\bm{B}_2\bm{B}_1$ and $\bm{R}=\bm{B}_2\bm{B}_1$. In a general setting, to ensure $\bm{R}\in\mathbb{R}^{d\times d}$ is dense, the factor matrices $\bm{B}_m,\ldots,\bm{B}_1$ must collectively generate a set of directed edges on a $d\times(m+1)$ grid, where each directed edge links two nodes in neighboring layers (\emph{i.e.}, consecutive columns), and every node from the first layer must have a pathway to every node in the final $(m+1)$-th layer.

The information transmission perspective offers valuable insights into the sparsity structure found in the factorization matrices of OFT. As depicted in Figure~\ref{fig:blk_diag}, the matrix $\bm{R}$ in the original OFT is characterized by a block-diagonal configuration. While this architecture serves to decrease the quantity of learnable parameters, it also precludes $\bm{R}$ from being a dense matrix. The objective is to synthesize a dense orthogonal matrix from $m$ sparse orthogonal matrices, achieving this with a minimal number of effective, trainable parameters. Within the information transmission paradigm, two core criteria guide us: \emph{(i)} \textbf{dense connectivity}, ensuring that every node at the initial layer can reach each node in the final layer through at least one path, and \emph{(ii)} \textbf{minimum free edges}, seeking the lowest total edge count possible without violating the orthogonality constraint. Orthogonality, in turn, imposes intricate restrictions on connections across adjacent layers. To illustrate, each $\bm{B}_i$ must be full-rank to meet the orthogonality requirement, necessitating $d$ connections that create a bijection between the nodes at the $i$-th and $(i=1)$-th levels. This establishes a lower bound of $d$ edges between neighboring levels (see the example in Figure~\ref{fig:info_trans}, which uses 4). These specific $d$ edges are essential for orthogonality and are excluded from our edge count since the corresponding elements are not tunable (\emph{e.g.}, in a $d\times d$ orthogonal matrix with $d$ non-zero components, the entries must be $\pm 1$). Because orthogonal matrices inherently require fewer parameters than unconstrained full matrices, enforcing orthogonality introduces extra dependencies among the edges. For instance, in the context of a $2\times 2$ orthogonal matrix, setting the $(1,1)$ element to zero automatically forces the $(2,2)$ element to zero as well (\emph{i.e.}, omitting one edge necessitates omitting another). Thus, depending on how edges are configured, the orthogonality condition might sometimes necessitate the addition or elimination of certain connections. By examining the support of the composed orthogonal matrix, the information transmission approach is particularly pertinent to OFT, as only the pattern of non-zero trainable entries in $\bm{R}$ is relevant, while their precise values are not.

A straightforward approach to forming dense connections between two hierarchical levels requires $\mathcal{O}(d^2)$ edges (that is, equivalent to one dense orthogonal matrix), amounting to $d^2-d$ edges—for instance, with $d=4$, this results in 12 edges. Figure~\ref{fig:info_trans} demonstrates a possible matrix factorization that instead utilizes just $10$ edges, which is fewer than the number needed for a full dense orthogonal matrix. This approach allows us to analyze the parameter-efficiency of OFT by using a graphical model, and it offers an intuitive way to devise alternative feasible matrix factorizations. Drawing on the structure of butterfly graphs from the Cooley-Tukey algorithm~\cite{cooley1965algorithm}, we note that these topologies are capable of connecting $d$ inputs to $d$ outputs using only $\mathcal{O}(d\log d)$ edges, providing dense connectivity with significantly fewer edges. As shown in Figure~\ref{fig:info_trans}, the example topology requires $10$ edges, while the butterfly network can achieve the same connectivity using just $8$ edges. We now present the butterfly structure and illustrate its benefits in terms of parameter efficiency.

\section{Orthogonal Parameterization by Butterfly Factorization}

The butterfly structure, originally utilized within the Cooley-Tukey algorithm for efficient computation of the fast Fourier transform, supports a paradigm where localized modifications in the frequency domain yield global alterations in the spatial domain. This property mirrors our context of information transmission, wherein each node at the initial layer is able to communicate with every node in the final layer. Furthermore, the butterfly configuration is widely adopted as a high-efficiency topology for computer networks~\cite{solihin2015fundamentals,li2011linear} to optimize data flow. For $k\geq 2$ where $k$ is a power of $2$, we define the \emph{butterfly factor} $\bm{B}^F(k)$ as follows:

\begin{equation}\label{eq:but_factor}
\begin{aligned}
    \bm{B}^F(k)=\begin{bmatrix}
    \text{diag}(\bm{d}_1) & \text{diag}(\bm{d}_2) \\
    \text{diag}(\bm{d}_3) & \text{diag}(\bm{d}_4)
    \end{bmatrix}\in\mathbb{R}^{k\times k},
\end{aligned}
\end{equation}

where each $\bm{d}_i\in\mathbb{R}^{\frac{k}{2}}, \forall i$ represents a vector. When $d=2^N$, we can then define the $d$-dimensional \emph{butterfly matrix} $\bm{B}(d)\in\mathbb{R}^{d\times d}$ recursively as

\begin{equation}
\begin{aligned}
    \!\!\bm{B}(d)=\tilde{\bm{B}}(d,d)\!\cdot\!\begin{bmatrix}
    \bm{B}_1(\frac{d}{2})\!\!\!\!\! & \bm{0} \\
    \bm{0}\!\!\!\!\! &\bm{B}_2(\frac{d}{2})
    \end{bmatrix}= \tilde{\bm{B}}(d,d)\tilde{\bm{B}}(d,\frac{d}{2})\cdots\tilde{\bm{B}}(d,2),
\end{aligned}
\end{equation}

Here, $\bm{B}_1(\frac{d}{2})$ and $\bm{B}_2(\frac{d}{2})$ represent butterfly matrices of dimension $\frac{d}{2}$. Next, we introduce the concept of a \emph{butterfly component}, denoted as $\thickmuskip=2mu \medmuskip=2mu \tilde{\bm{B}}(d,k)=\text{diag}(\bm{B}^F_1(k),\cdots,\bm{B}^F_{d/k}(k))$, which is a block-diagonal matrix of size $d \times d$ with blocks of size $k$. Here, each $\bm{B}^F_i(k)$, for all $i$, refers to a butterfly factor as specified in Equation~\ref{eq:but_factor}. This formulation allows us to utilize butterfly matrices for the parameterization of orthogonal matrices by ensuring that every multiplicative component $\tilde{\bm{B}}(d,k)$ of the butterfly matrix $\bm{B}(d)$ is orthogonal. We begin by analyzing $\tilde{\bm{B}}(d,2)$, a block-diagonal matrix with block size 2. It is straightforward to guarantee the orthogonality of $\tilde{\bm{B}}(d,2)$ using the Cayley transform, or equivalently through parameterizing each block as a $2$-dimensional rotation, as described in \cite{liu2021orthogonal,qiu2023controlling}. The sparsity structure of each butterfly component corresponds to a permuted arrangement of the non-zero entries in $\tilde{\bm{B}}(d,2)$, allowing every butterfly component to also be easily parameterized as an orthogonal matrix. This results in an effective parameterization strategy for orthogonal matrices based on aggregating multiple $2 \times 2$ orthogonal matrices. Following the approach of \cite{chen2022pixelated}, we extend the concept of butterfly matrices to define a block butterfly component, $\tilde{\bm{B}}^b(d,k)$, where every element in $\bm{d}_i$, for any $i$, is replaced with a $b \times b$ matrix. To ensure that the block butterfly component $\tilde{\bm{B}}^b(d,2)$ remains orthogonal, we parameterize each $2b \times 2b$ block matrix as orthogonal. For other butterfly components $\tilde{\bm{B}}^b(d,k)$ with $k>2$, their non-zero configurations result from block-wise permutations of those in $\tilde{\bm{B}}^b(d,2)$, which means orthogonality can similarly be maintained. Integrating these elements, the forward computation in BOFT is as follows:

\begin{equation}
    \begin{aligned}\nonumber
    \bm{z} = \big(\bm{R}(m, b) \cdot \bm{W}^0\big)^\top \bm{x}, \qquad \text{s.t.}~\bigg{\{} \bm{R}(m, b) = \prod_{i=1}^m \tilde{\bm{B}}^b_{(d, i)} ~~\&~~ \underbrace{\big(\tilde{\bm{B}}^b_{(d, j)}\big)^\top \tilde{\bm{B}}^b_{(d, j)} = \tilde{\bm{B}}^b_{(d, j)} \big(\tilde{\bm{B}}^b_{(d, j)}\big)^\top = \bm{I}_d}_{\forall j\in [1, m]} \bigg{\}},
    \end{aligned}
\end{equation}

For notation brevity, $\tilde{\bm{B}}^b(d, 2^{m-i+1})$ is written as $\tilde{\bm{B}}^b_{(d, i)}$, and $\bm{I}_d$ refers to the $d\times d$ identity matrix. The orthogonal matrix $\bm{R}(m, b)\in\mathbb{R}^{d\times d}$ is constructed by multiplying several orthogonal butterfly matrices together. BOFT using $\bm{R}(m, \frac{b}{2})$ is simply denoted as BOFT($m$, $b$) for $b\geq 2$. Specializing to the case $m=1$, BOFT($1$, $b$) becomes equivalent to the block-diagonal OFT~\cite{qiu2023controlling} with block size $b$. In the case of BOFT($1$, $d$), the scheme recovers the standard OFT~\cite{qiu2023controlling}, corresponding to an unconstrained orthogonal matrix. Furthermore, setting $m = \log \frac{2d}{b}$ in BOFT($m$, $b$) employs the block butterfly matrix $\bm{B}^{\frac{b}{2}}(d)$ as $\bm{R}$, resulting in a dense orthogonal transformation. Generally, BOFT($m$, $b$) utilizes $\frac{1}{2}(b-1)dm$ effective parameters for tuning a linear layer with dimensions $d\times n$. When using the butterfly architecture, i.e., $m = \log d$, $b = 2$, BOFT has parameter complexity $\mathcal{O}(d\log d)$. In contrast, classic OFT employing a dense orthogonal matrix needs $\mathcal{O}(d^2)$ parameters, while block-diagonal OFT with $r$ blocks requires $\mathcal{O}(bd)$. Thus, the conventional OFT must set $b = d$ to create a fully dense orthogonal matrix, whereas BOFT achieves a dense matrix for any block size $b$.

\textbf{Identity initialization in BOFT}. Commonly, parameter-efficient finetuning initializes from the exact parameters of the pretrained model to ensure the finetuned network remains close to the original. LoRA, for example, applies zero initialization for its low-rank components. In BOFT, initialization sets every butterfly matrix component to be the identity (\emph{i.e.}, initializing the associated skew-symmetric matrices as zeros in the Cayley parameterization).

\textbf{Multiplicative Dropout adaptation for BOFT}. While LoRA~\cite{hulora2022} incorporates a Dropout layer for its low-rank weight update to curb overfitting and conventional Dropout~\cite{srivastava2014dropout} is effective for LoRA, it is incompatible with BOFT due to the multiplicative update mechanism. To resolve this, we design a multiplicative Dropout tailored to BOFT. Since $\bm{R}(m, b)$ comprises $m$ orthogonal butterfly matrices that can be reordered into $2b\times 2b$ block-diagonal orthogonal forms, this Dropout variant proceeds by randomly selecting $p_1$ proportion of the butterfly factors and $p_2$ fraction of diagonal blocks within each component, and substituting the chosen blocks with identity matrices.

\section{Intriguing Insights and Discussions}

\textbf{Expressivity of BOFT}. The butterfly architecture combined with permutations can exactly reconstruct a range of established fast linear transforms~\cite{dao2019learning,dao2019kaleidoscope} (\emph{e.g.}, fast Fourier transform, Hadamard transform). However, the capacity of our orthogonal butterfly matrix to approximate arbitrary orthogonal matrices has yet to be thoroughly evaluated. To investigate this, we performed simulations aiming to approximate a random, dense orthogonal matrix~\cite{anderson1987generation} of dimension $512\times 512$. Figure~\ref{fig:para_eff} displays results averaged across 10 independent random seeds. On these plots, the y-axis indicates the approximation error, while the x-axis represents the effective number of trainable parameters. Each colored curve represents a BOFT configuration with a fixed block size; the point at the extreme left denotes BOFT($1$,$b$) (\emph{i.e.}, the standard block-diagonal OFT for block size $b$). Across these experiments, BOFT consistently achieves greater parameter efficiency compared to OFT. For instance, BOFT($9$,$2$) demonstrates superior expressivity compared to BOFT($1$,$16$), while utilizing significantly fewer parameters. In general, BOFT achieves higher parameter efficiency with smaller block sizes ($b$) and larger butterfly factors ($m$); for example, BOFT($6$,$4$) requires far fewer parameters but yields an approximation error comparable to BOFT($2$,$16$). In summary, the butterfly matrix structure defines a more organized subset within the orthogonal group compared to block-diagonal constructions, thus allowing BOFT to be provably more expressive than OFT when using identical block sizes.

\begin{theorem}[Expressivity of BOFT]\label{thm:exp_boft}
    BOFT surpasses OFT in expressivity when the block sizes are equal. Any orthogonal matrix of size $d$ can be approximated by a product of butterfly matrices: $\bm{B}_{d-1,1}(d)\bm{B}^\top_{d-1,2}(d)\cdots\bm{B}_{1,1}(d)\bm{B}^\top_{1,2}(d)$, where each $\bm{B}_{i,j}(d),\forall i,\forall j$ is a butterfly matrix.
\end{theorem}

The result in Theorem~\ref{thm:exp_boft} motivates an extended form for BOFT, where the final orthogonal matrix is constructed as $\thickmuskip=2mu \medmuskip=2mu \bm{R}^G(m_1,b_1,m_2,b_2,l)=\bm{R}_{l,1}(m_1,b_1)\bm{R}_{l,2}^T(m_2,b_2)\cdots\bm{R}_{1,1}(m_1,b_1)\bm{R}_{1,2}^T(m_2,b_2)$. Here, $\bm{R}_{i,j}^T(m,b)$ represents the orthogonal matrix used in BOFT. Setting $m_1=m_2=\log d$, $b_1=b_2=2$, and $l=d-1$, allows $\bm{R}^G(m_1,b_1,m_2,b_2,l)$ to cover the complete orthogonal group, a structure referred to as the kaleidoscope hierarchy~\cite{dao2019kaleidoscope}. Nevertheless, it is important to note that higher expressivity does not guarantee superior outcomes during finetuning. For instance, fully parameterized finetuning, despite its theoretical universality, often performs poorly in practice. Therefore, navigating the balance between expressivity and regularity is critical for robust finetuning. By infusing structural priors, BOFT expands the available parameter space for finetuning, thus facilitating a more effective trade-off between expressivity and regularization.

\textbf{Spectral properties}. Orthogonal finetuning typically results in superior spectral characteristics compared to LoRA, as it maintains the spectral norm of the original pretrained weight matrix $\bm{W}^0$ exactly. To understand this, consider the singular value decomposition of $\bm{W}^0$: $\bm{W}^0=\bm{U}\bm{\Sigma}\bm{V}^\top$, where $\bm{U}$ and $\bm{V}$ are orthogonal matrices and $\bm{\Sigma}$ is a diagonal matrix containing the singular values. In both OFT and BOFT, an orthogonal transformation $\bm{R}$ is applied to the left, yielding finetuned weights of the form $\bm{R}\bm{U}\bm{\Sigma}\bm{V}^\top$. This operation leaves the largest singular value—and thus the spectral norm of $\bm{W}^0$—unaltered. Maintaining this property has been demonstrated to substantially promote both training stability and model generalization~\cite{miyato2018spectral,yoshida2017spectral}. Further mathematical details are provided in Appendix~\ref{app:sn_property}.

\textbf{Orthogonal finetuning as learning bilinear similarity}. The BOFT approach can be formulated in terms of learning a bilinear similarity function $\bm{w}^0_i\bm{R}\bm{x}$, where $\bm{w}^0_i$ denotes the $i$-th neuron (or column) of $\bm{W}^0$. In this context, BOFT can be interpreted as optimizing the bilinear similarity matrix $\bm{R}$, with the additional constraint that $\bm{R}$ must be orthogonal. This perspective naturally ties BOFT to concepts in distance metric learning~\cite{xing2002distance} and bilinear forms~\cite{roman2005advanced}.

\textbf{Inductive bias and generalization in BOFT}. The matrix $\bm{R}(m,b)$ in BOFT frequently embodies a structured portion of the orthogonal group, imposing specific restrictions on the space of possible hypotheses. This structural constraint instills an inductive bias. We contend that the inductive bias emerging from the butterfly factorization is advantageous for generalization, as it echoes the structured patterns found in well-known linear operators~\cite{dao2019kaleidoscope}, such as the discrete Fourier, sine/cosine, and Hadamard transforms. Additionally, the sparse factorization property within BOFT might induce further implicit inductive biases~\cite{gunasekar2017implicit,li2018algorithmic,liu2019neural}.

\textbf{Contrast with butterfly-based sparse training}. A number of studies~\cite{dao2019learning,dao2019kaleidoscope,chen2022pixelated,dao2022monarch} have explored sparse training utilizing butterfly parameterizations. These approaches primarily involve directly reparameterizing neural network weight matrices through butterfly structures and subsequently training the networks from initialization. Notably, \cite{dao2022monarch} investigates finetuning by projecting pretrained weights onto a modified butterfly matrix and updating the projected parameters for downstream applications. In contrast, BOFT introduces a novel finetuning approach, employing transformations of the weights using weight matrices shared across layers.

\input{rebuttal-results}

\section{Concluding Remarks and Limitations}

This work introduces Orthogonal Butterfly, a general parameter-efficient finetuning technique for foundation models grounded in the butterfly framework. Our central idea for improving parameter efficiency lies in representing a dense orthogonal matrix as a product of several sparse orthogonal matrices. To streamline the discovery of appropriate matrix factorizations, we present a graphical information transmission perspective. Using this framework, we reveal that the butterfly structure is particularly suitable for constructing sparse orthogonal matrix decompositions that meet our goals. We empirically establish BOFT's strong performance when finetuning a range of models, including large language models, substantial vision architectures, and text-to-image generative systems. Experimental results further affirm BOFT's advantages as a broadly applicable finetuning scheme.

Nevertheless, BOFT is not without its drawbacks. The construction of the final orthogonal matrix via multiplication of multiple orthogonal matrices introduces a slightly increased training runtime relative to OFT, and optimizing BOFT’s training efficiency remains unresolved. On the upside, once finetuning is complete, the orthogonal matrices acquired by BOFT may be merged into the pretrained network without incurring additional inference time. Furthermore, it is still uncertain whether the butterfly topology is the optimal structure for information transfer. Our information transmission approach opens possibilities for leveraging insights from computer networking research, which systematically studies network topology efficiency for data flow. We anticipate that future exploration of alternative network structures could yield even more efficient designs for synthesizing dense orthogonal matrices.

\end{document}